% Conduct
\subsection*{Professional Conduct}
Cheating and plagiarism are serious academic offenses that Colleges and Universities deal with firmly. Academic integrity presumes that students are honest in all academic work. Cheating is the failure to give credit for work not done independently (i.e., submitting a paper written by someone other than yourself), unauthorized communication dABCng an examination, or claiming credit for work not done (i.e., falsifying information). Plagiarism is the failure to give credit for another person’s written or oral statement, thereby falsely presuming that such work is originally and solely your own. \\

Students are expected to be honest in all academic work. A student’s name on any written work, quiz, or exam shall be regarded as assurance that the work results from the student’s independent thought and study. Work should be stated in the student’s words and correctly attributed to its source. Students should know how to quote, paraphrase, summarize, cite, and reference the work of others with integrity. The following are examples of academic dishonesty.
\begin{outline}
    \1	Using material, directly or paraphrasing, from published sources (print or electronic) without appropriate citation;
    \1	Claiming disproportionate credit for work not done independently;
    \1	Unauthorized possession or access to exams;
    \1	Unauthorized communication during exams;
    \1	Unauthorized use of another’s work or preparing work for another student;
    \1	Taking an exam for another student;
    \1	Altering or attempting to alter grades;
    \1	The use of notes or electronic devices to gain an unauthorized advantage during exams;
    \1	Fabricating or falsifying facts, data, or references;
    \1	Facilitating or aiding another’s academic dishonesty;
    \1	Submitting the same paper for more than one course without prior approval from the Instructor.
\end{outline}

\\

\vspace{2ex}\hrule\vspace{2ex}